\scanpage{75}
\uline{Catégories} \uline{\unsure{hopolynômiales|polynômiales|polynomiales}}

Une \uline{catégorie polynômiale} est une catégorie $P$,
munie d'une sous-catégorie (pas pleine en général)
$P_1$, avec les axiomes suivants
\margin{Notation $P(X,Y)$ pour $\operatorname{Hom}_P(X,Y)$}

a) $P$ admet des produits finis, \deleted{et l'objet final est}

b) $\operatorname{Ob}P_1=\operatorname{Ob}P$, \deleted{les produits et pour deux objets
dans $P$, leur produit dans $P$ est un produit
dans $P_1$ (i.e. si $X,Y,Z\in\operatorname{Ob}P_1$, \deleted{$X,Y\in\operatorname{Ob}P$} \unsure{l'ensemble de}
\inserted{(i.e si $X,Y,Z\in\operatorname{Ob}P_1$, une} flèche $Z\to X\times Y$ de $P$ est dans $P_1$ ssi
$(\unsure{u,v|v,w}) : X\to Y\times Z$ est dans $P_1$ ssi} (donc $P_1$ est décrite
par un ens\unsure{.} $\mathrm{Fl}_1\subset\mathrm{Fl}(P)=\mathrm{Fl}$ contenant les identités,
et stable par composition.) On veut de plus que
\deleted{$\mathrm{Fl}_1$ contienne les isos, et que} pour $X,Y\in\operatorname{Ob}P$,
leur produit $X\times Y$ soit un produit ds $P_1$
(i.e. une flèche $u=(\unsure{u',u''|u,v|v,w}) : Z\to X\times Y$ est
``de degré $1$'' ssi ses composantes $X,Y$ le sont),
\inserted{enfin si $0$ est l'objet final, la fl $X\to0$ est de degré $1$ (donc $0$ objet final de $P_1$)}

c) $P_1$ est une catégorie additive

Cela implique que \deleted{pour tt objet l'on a
un hom de monoïdes,} les $\operatorname{Hom}_P(X,Y)$
\deleted{\[
\mathbb Z\longrightarrow\operatorname{End}(\mathrm{id}_{P_1})\longleftarrow\operatorname{End}(\mathrm{id}_P)
\]
\unsure{de plus} pr tt $\deleted{\unsure{A}}X,Y\in\operatorname{Ob}P$, \unsure{ou un} $X$}
\deleted{i.e pour tt $X$.}
sont des groupes abéliens (car les objets $\deleted{X}$\inserted{de $P$} sont
des objets groupes abéliens de $P$), dépendent
fonctorielt de $X$, mais pas de $Y$
(i.e. on a
\[
\boxed{(f+g)\circ h=fh+g\circ h}
\]
mais \deleted{\unsure{non en général}}\inserted{on \unsure{peut|va} avoir,} éventuellement
\[
f\circ(g+h)\ne f\circ g+f\circ h\quad !
\]

Il en résulte aussi que $\mathbb Z^\times$ opère sur $P$
i.e on a un homomorphisme multiplicatif
\[
\mathbb Z^\times\longrightarrow\operatorname{End}(P_1)\longleftarrow\operatorname{End}(P)
\]
qui à $n\in\mathbb Z$ et $X\in\operatorname{Ob}P_1=\operatorname{Ob}P$ associe $n\,\mathrm{id}_X$.
Ainsi $\mathbb Z$ opère dans $\deleted{\operatorname{Hom}_{\unsure{P}}}P(X,Y)$ par
$\mu_n(f)=f\circ(n\,\mathrm{id}_Y)$
% archive sheet number: 74
